%%%%%%%%%%%%%%%%%%%%%%%%%%%%%%%%%%%%%%%%%%%%%%%%%%
%% 第十二章 素数定理（证明提要）
%% 详细估计见附录 E（\ref{app:pnt_estimates}）
%%%%%%%%%%%%%%%%%%%%%%%%%%%%%%%%%%%%%%%%%%%%%%%%%%

\chapter{素数定理}\label{ch:prime_number_theorem}

\section*{引言}\label{sec:pnt_ch_intro}

第~\ref{ch:prime_distribution}~、\ref{ch:elementary_number_theoretic_functions}~章给出了古典经验公式
$x/\log x$、$\Li(x)$ 与素数定理的\textbf{陈述}；
第~\ref{ch:riemann_zeta_continuation}~章建立了 $\zeta$ 的延拓与函数方程；
第~\ref{ch:xi_function}~章给出 $\xi$ 的 Hadamard 乘积与 $\zeta'/\zeta$ 的展开；
第~\ref{ch:zero_distribution}~章说明非平凡零点位于开临界带，并强调
$\zeta(1+it)\neq 0$（$t\neq 0$）在历史上与素数定理的联系。

本章给出素数定理的\textbf{证明提要}。安排如下：
\begin{enumerate}
  \item \textbf{目标与等价}：将 $\pi\sim\Li\sim x/\log x$ 化为 $\psi\sim x$
    （第~\ref{ch:elementary_number_theoretic_functions}~章已证的初等等价）；
  \item \textbf{解析输入}：汇总 $-\zeta'/\zeta$、延拓、Hadamard 展开，
    并陈述关键输入 $\zeta(1+it)\neq 0$（细节见附录~\ref{app:pnt_estimates}）；
  \item \textbf{Perron 与主项}：由围道积分取出 $s=1$ 处留数 $x$，说明为何需要线上无零点，
    并简要说明其余贡献为 $o(x)$；
  \item \textbf{结论}：由初等等价从 $\psi\sim x$ 得 $\pi\sim\Li$，并与下一章显式公式划界
    （本章只管主项合法性，误差结构见第~\ref{ch:riemann_explicit_formula}~章）。
\end{enumerate}
完整的三角不等式、零点自由区与围道边估计见\textbf{附录~\ref{app:pnt_estimates}}；
本章不重复附录中的长计算。

逻辑流程见图~\ref{fig:pnt-flow}。

\begin{figure}[htbp]
\centering
\begin{tikzpicture}[
  font=\small,
  >=Stealth,
  node distance=11mm and 18mm,
  main/.style={
    draw=MainBlue!65!black,
    fill=LightBlue!35,
    rounded corners=3pt,
    align=center,
    inner sep=7pt,
    text width=5.6cm,
    minimum height=1.15cm,
  },
  side/.style={
    draw=Gold!75!black,
    fill=LightGold!55,
    rounded corners=3pt,
    align=center,
    inner sep=6pt,
    text width=3.2cm,
    minimum height=1.05cm,
  },
  res/.style={
    draw=DarkGreen!65!black,
    fill=LightGreen!40,
    rounded corners=3pt,
    align=center,
    inner sep=7pt,
    text width=5.6cm,
    minimum height=1.15cm,
  },
  arr/.style={->, very thick, MainBlue!75!black},
  sarr/.style={->, thick, Gold!70!black},
  lab/.style={font=\footnotesize\sffamily, text=MainBlue!65!black},
]
  % -------- 主流程（纵向中轴）--------
  \node[main] (N1) {%
    \textbf{1.\ 解析输入}\\[1pt]
    延拓与 $-\zeta'/\zeta$\\
    {\footnotesize（第 5、9、10 章）}
  };
  \node[main, below=of N1] (N2) {%
    \textbf{2.\ Perron 积分}\\[1pt]
    $\psi(x)=\frac1{2\pi i}\int_{(c)}\Bigl(-\frac{\zeta'}{\zeta}\Bigr)\frac{x^s}{s}\,\mathrm{d}s$\\
    {\footnotesize（第 4、6 章；\eqref{eq:psi_perron}）}
  };
  \node[main, below=of N2] (N3) {%
    \textbf{3.\ 主项留数}\\[1pt]
    $\operatorname{Res}_{s=1}=x$
  };
  \node[res, below=16mm of N3] (N4) {%
    \textbf{4.\ $\psi(x)\sim x$}\\[1pt]
    即 $\psi(x)=x+o(x)$
  };
  % N4–N5 加大竖距，供「余项控制」落在间隙、「初等等价」贴 N5
  \node[res, below=32mm of N4] (N5) {%
    \textbf{5.\ 素数定理}\\[1pt]
    $\pi\sim\Li\sim x/\log x$
  };
  \node[main, below=of N5] (N6) {%
    \textbf{6.\ 下一章}\\[1pt]
    显式公式 $\psi_0=x-\sum_{\rho}x^{\rho}/\rho+\cdots$\\
    {\footnotesize（同一主项 $x$，展开误差结构）}
  };

  % -------- 侧翼（主轴左侧）--------
  % S1：N4 上方外侧；S2：N4–N5 正中外侧；S3：N5 外侧
  \node[side] (S1) at ([xshift=-4.9cm, yshift=1.85cm]N4.west) {%
    \textbf{$\operatorname{Re}s=1$ 上无零点}\\[1pt]
    $\zeta(1+it)\neq 0$\\
    {\footnotesize（$t\neq 0$；附录~\ref{app:pnt_estimates}）}
  };
  \node[side] (S2) at ($($(N4.west)!0.38!(N5.west)$)+(-4.9cm,0)$) {%
    \textbf{余项控制}\\[1pt]
    自由区\,/\,边估计\\
    $\Rightarrow$ 余项 $o(x)$\\
    {\footnotesize（附录~\ref{app:pnt_estimates}）}
  };
  % 略低于 N5 中心，与 S2 再拉开
  \node[side] (S3) at ([xshift=-4.9cm, yshift=-1.15cm]N5.west) {%
    \textbf{初等等价}\\[1pt]
    $\psi\sim x\Leftrightarrow\pi\sim x/\log x$\\
    {\footnotesize（第 2 章；$\Li\sim x/\log x$）}
  };

  % -------- 主轴纵向箭头 --------
  \draw[arr] (N1) -- (N2);
  \draw[arr] (N2) -- (N3);
  \draw[arr] (N3) -- (N4);
  \draw[arr] (N4) -- (N5);
  \draw[arr] (N5) -- (N6);

  % -------- 侧翼从边接入主轴 --------
  \draw[sarr] (S1.east) -- (N4.north west);
  \draw[sarr] (S2.east) -- (N4.south west); % 余项并入 $\psi\sim x$
  \draw[sarr] (S3.east) -- (N5.south west); % 初等等价并入素数定理

  % -------- 旁注（主轴右侧，略外移）--------
  \node[lab, right=5mm of N3] {主项来源};
  \node[lab, right=5mm of N4] {主项$+$无零点$+$余项};
  \node[lab, right=5mm of N5] {化为经典形式};
\end{tikzpicture}
\caption{本章证明提要流程图（主轴自上而下）。
中轴：$1\to 6$ 为证明主流程；
左侧金框经侧边接入：$\zeta(1+it)\neq 0$ 与自由区\,/\,边估计保证 $\psi=x+o(x)$，
初等等价将 $\psi\sim x$ 化为 $\pi\sim\Li$。
$3$--$4$--$1$ 与边估计细节见附录~\ref{app:pnt_estimates}。}
\label{fig:pnt-flow}
\end{figure}

%===========================================================
\section{目标与等价形式}\label{sec:pnt_goal}
%===========================================================

\begin{theorem}{素数定理}{}
\label{thm:pnt-main}
当 $x\to\infty$ 时，
\[
  \pi(x)\sim\Li(x)\sim\frac{x}{\log x}.
\]
\end{theorem}

第~\ref{ch:elementary_number_theoretic_functions}~章已证（在初等范围内）下列渐近等价：
\begin{equation}
  \pi(x)\sim\frac{x}{\log x}
  \;\Longleftrightarrow\;
  \theta(x)\sim x
  \;\Longleftrightarrow\;
  \psi(x)\sim x.
  \label{eq:pnt-equiv-ch11}
\end{equation}
又 $\Li(x)\sim x/\log x$（第~\ref{ch:prime_distribution}~章分部积分）。
因此只需证明解析形式
\begin{equation}
  \boxed{\psi(x)\sim x}\,,
  \label{eq:psi-sim-x}
\end{equation}
即可得到定理~\ref{thm:pnt-main}。

\begin{remark}{与显式公式的分工}{}
第~\ref{ch:riemann_explicit_formula}~章的显式公式给出
\[
  \psi_0(x)=x-\sum_{\rho}\frac{x^{\rho}}{\rho}-\log(2\pi)-\frac12\log(1-x^{-2})
\]
（跳跃点取平均）。本章只证明\textbf{主项} $\psi\sim x$；
零点级数控制的是误差 $\psi(x)-x$，其阶与 RH 相关，不属于素数定理本身。
注意逻辑方向：显式公式是恒等式，把 $\psi-x$ 写成零点级数；素数定理是估计，要求该级数为 $o(x)$。
前者\textbf{并不蕴含}后者——由公式过渡到 $\psi\sim x$，仍须 $\zeta(1+it)\neq 0$ 及相应余项控制（见下文与附录~\ref{app:pnt_estimates}）。
\end{remark}

%===========================================================
\section{解析输入}\label{sec:pnt_inputs}
%===========================================================

\subsection{Dirichlet 级数}
对 $\operatorname{Re}s>1$（第~\ref{ch:riemann_zeta_intro}~章），
\begin{equation}
  -\frac{\zeta'(s)}{\zeta(s)}
  =\sum_{n=1}^{\infty}\frac{\Lambda(n)}{n^{s}}.
  \label{eq:pnt-logderiv}
\end{equation}
$\psi(x)=\sum_{n\le x}\Lambda(n)$ 正是其系数的部分和。

\subsection{延拓与极点}
由第~\ref{ch:riemann_zeta_continuation}~章，$\zeta$ 亚纯延拓至 $\mathbb{C}$，仅在 $s=1$ 有单极点，
$\operatorname{Res}_{s=1}\zeta=1$。
因而 $-\zeta'/\zeta$ 在 $s=1$ 有单极点，留数为 $1$
（与 $\zeta(s)\sim 1/(s-1)$ 一致）。

\subsection{Hadamard 与对数导数}
第~\ref{ch:xi_function}~章给出
\begin{equation}
  \frac{\zeta'(s)}{\zeta(s)}
  =b_0+\sum_{\rho}\Bigl(\frac{1}{s-\rho}+\frac{1}{\rho}\Bigr)
  -\frac{1}{s}-\frac{1}{s-1}
  +\frac12\log\pi
  -\frac12\frac{\Gamma'(s/2)}{\Gamma(s/2)},
  \label{eq:pnt-zeta-from-xi}
\end{equation}
（在对称截断约定下）。由此可读出：$-\zeta'/\zeta$ 的极点来自 $s=1$、非平凡零点、平凡零点与 $s=0$ 等。

\subsection{关键输入：\texorpdfstring{$\operatorname{Re}s=1$}{Re s=1} 上无零点}
\begin{theorem}{直线 $\operatorname{Re}s=1$ 上无零点}{}
\label{thm:no-zero-on-one}
对一切实数 $t\neq 0$，
\[
  \zeta(1+it)\neq 0.
\]
\end{theorem}

第~\ref{ch:riemann_zeta_intro}~章已证 $\operatorname{Re}s>1$ 时 $\zeta\neq 0$；
边界 $s=1+it$ 不能再用 Euler 乘积绝对收敛直接得到。
Hadamard 与 de la Vallée Poussin（1896）给出的经典路线使用
$3$--$4$--$1$ 型三角不等式：对 $\sigma>1$，
\begin{equation}
  \zeta(\sigma)^{3}\,
  \bigl|\zeta(\sigma+it)\bigr|^{4}\,
  \bigl|\zeta(\sigma+2it)\bigr|
  \ge 1,
  \label{eq:pnt-341}
\end{equation}
再令 $\sigma\to 1^+$ 并排除 $\zeta(1+it)=0$ 的可能性。
\textbf{完整推导见附录~\ref{app:pnt_estimates} 第~\ref{sec:app-no-zero-Re1}~节}；
第~\ref{ch:zero_distribution}~章已将其作为临界带右边界的陈述使用。

\begin{remark}{与黎曼猜想（RH）的区别}{}
定理~\ref{thm:no-zero-on-one} 只禁止零点落在 $\operatorname{Re}s=1$ 上，
\textbf{不}要求 $\operatorname{Re}\rho=1/2$。
素数定理弱于 RH；后者控制的是误差阶。
\end{remark}

%===========================================================
\section{从 Perron 积分到主项}\label{sec:pnt_perron_outline}
%===========================================================

\subsection{Perron 公式}
出发点取第~\ref{ch:discrete_continuous}~章已建立的~$\psi$ 的 Perron 表示~\eqref{eq:psi_perron}
（$c>1$；非整数 $x$ 时取 $\psi(x)$，跳跃点取平均 $\psi_0$）：
\[
  \psi(x)
  =\frac{1}{2\pi i}
  \int_{c-i\infty}^{c+i\infty}
  \Biggl(-\frac{\zeta'(s)}{\zeta(s)}\Biggr)
  \frac{x^{s}}{s}\,\mathrm{d}s.
\]
（一般 Perron 见第~\ref{ch:integral_transforms}~章；截断与误差见附录~\ref{app:pnt_estimates}。）
直观上：把垂直线向左平移，用留数收集 $-\zeta'/\zeta\cdot x^{s}/s$ 的极点贡献。

\subsection{主项留数}
在 $s=1$ 处，$-\zeta'/\zeta$ 的留数为 $1$，$x^{s}/s$ 全纯且取值 $x$，故
\begin{equation}
  \operatorname{Res}_{s=1}
  \Biggl[\Biggl(-\frac{\zeta'(s)}{\zeta(s)}\Biggr)\frac{x^{s}}{s}\Biggr]
  =x.
  \label{eq:pnt-residue-one}
\end{equation}
这正是 $\psi(x)\sim x$ 的主项来源，与第~\ref{ch:riemann_explicit_formula}~章显式公式中的 $x$ 相同。

\subsection{为何需要 $\zeta(1+it)\neq 0$}
若存在 $t_0\neq 0$ 使 $\zeta(1+it_0)=0$，则 $-\zeta'/\zeta$ 在 $s=1+it_0$ 有极点，
围道平移时多出阶约 $x^{1+it_0}/(1+it_0)$ 的项，其模为 $x$，
足以破坏 $\psi(x)-x=o(x)$。
因此定理~\ref{thm:no-zero-on-one} 是保证主项 $x$ 不被同量级极点破坏的\textbf{必要条件}。

\subsection{其余贡献的量级（提要）}
向左平移并取极限时，还需控制：
\begin{enumerate}
  \item 非平凡零点贡献 $\sum_{\rho}x^{\rho}/\rho$：在仅知 $\operatorname{Re}\rho<1$ 时，
    不足以单独得到 $o(x)$，而需零点自由区或等价的密度/增长估计
    （Hadamard--de la Vallée Poussin 零点自由区，见附录~\ref{app:pnt_estimates}）；
  \item 平凡零点与 $s=0$ 等：贡献 $O(1)$ 或更小；
  \item 水平边与左竖直边：在截断高度 $T\to\infty$ 时趋于 $0$
    （估计见附录及第~\ref{ch:integral_transforms}~章）。
\end{enumerate}
综合后得 $\psi(x)=x+o(x)$，即~\eqref{eq:psi-sim-x}。

\begin{remark}{两种现代写法}{}
经典路线强调零点自由区 $+$ 显式公式截断；
亦可采用 Wiener--Ikehara 或 Newman 的 Tauberian 路线
（第~\ref{ch:integral_transforms}~章曾提及），把
$-\zeta'/\zeta-1/(s-1)$ 在 $\operatorname{Re}s\ge 1$ 的性质直接化为 $\psi(x)\sim x$。
关键条件同样是 $\zeta(1+it)\neq 0$。细节见附录~\ref{app:pnt_estimates} 第~\ref{sec:app-tauberian}~节。
\end{remark}

%===========================================================
\section{结论：从 \texorpdfstring{$\psi\sim x$}{psi ~ x} 到 \texorpdfstring{$\pi\sim\Li$}{pi ~ Li}}\label{sec:pnt_conclude}
%===========================================================

由~\eqref{eq:psi-sim-x} 与~\eqref{eq:pnt-equiv-ch11}，
\[
  \pi(x)\sim\frac{x}{\log x}.
\]
又 $\Li(x)\sim x/\log x$（第~\ref{ch:prime_distribution}~章），故
\[
  \pi(x)\sim\Li(x)\sim\frac{x}{\log x},
\]
即定理~\ref{thm:pnt-main}。

这完成了第~\ref{ch:prime_distribution}~章预告的「主项合法性」：
古典 $\Li$ 与 $x/\log x$ 原为经验公式，素数定理则在解析上证明它们是 $\pi$ 的正确渐近主部。
误差 $\pi-\Li$ 的形状由下一章显式公式中的零点级数描述；
其最强常见阶与 RH 等价（von Koch），属于更强命题。

\begin{center}
\fbox{\parbox{0.92\textwidth}{
\centering
\vspace{0.2em}
\textbf{第十二章要点}\\[6pt]
\renewcommand{\arraystretch}{1.45}
\begin{tabular}{lp{9.0cm}}
\hline
目标 & $\psi\sim x\;\Rightarrow\;\pi\sim\Li\sim x/\log x$ \\
关键输入 & $\zeta(1+it)\neq 0$（$t\neq 0$）；延拓；$-\zeta'/\zeta$ \\
主项 & $\operatorname{Res}_{s=1}(-\zeta'/\zeta\cdot x^{s}/s)=x$ \\
详细估计 & 附录~\ref{app:pnt_estimates} \\
下文 & 第~\ref{ch:logint_Ei}~章（$\Li$/$\li$/$\Ei$）；
  第~\ref{ch:riemann_explicit_formula}~章：主项 $+$ 零点误差 \\
\hline
\end{tabular}
\vspace{0.2em}
}}
\end{center}

\section*{本章小结与后续展望}
\addcontentsline{toc}{section}{本章小结与后续展望}

本章在第~\ref{ch:zero_distribution}~章之后给出素数定理的证明提要：
以 $\zeta(1+it)\neq 0$ 与 Perron 积分为核心，说明 $\psi\sim x$ 如何导出 $\pi\sim\Li$。
完整的 $3$--$4$--$1$ 证明、零点自由区与边估计见附录~\ref{app:pnt_estimates}。

下一章推导黎曼显式公式：在已合法的主项之上展开零点贡献，
给出 $\psi$（及 $J$、$\pi$）与非平凡零点之间的精确对应；
素数定理解决的是主项主导，显式公式给出的是完整振荡结构。
